\chapter{จุลินทรีย์ในอุตสาหกรรมเกษตรแปรรูปและเทคโนโลยีชีวภาพ}
\label{chap:ch07}

% ==========================================================
% แผนบริหารการสอนประจำบทที่ 7
% ==========================================================
\section*{แผนบริหารการสอนประจำบทที่ 7}
\addcontentsline{toc}{section}{แผนบริหารการสอนประจำบทที่ 7}

\noindent\textbf{หัวข้อเนื้อหาประจำบท}
\begin{enumerate}
    \item ขอบข่ายของเทคโนโลยีชีวภาพจุลินทรีย์ในอุตสาหกรรมเกษตร
    \item จุลชีววิทยาของการหมักพืชอาหารสัตว์ (Silage Fermentation)
    \item การประยุกต์ใช้โพรไบโอติก (Probiotics) และพรีไบโอติกในการปศุสัตว์
    \item ชีววิทยาของการเพาะเห็ดเศรษฐกิจและการย่อยสลายเศษวัสดุเหลือทิ้ง
    \item การผลิตโปรตีนเซลล์เดียว (Single-Cell Protein: SCP) และสารชีวภัณฑ์มูลค่าสูง
\end{enumerate}

\noindent\textbf{วัตถุประสงค์เชิงพฤติกรรม}
\begin{enumerate}
    \item อธิบายกระบวนการหมักพืชอาหารสัตว์และบทบาทของแบคทีเรียกรดแลกติกได้
    \item วิเคราะห์ปัจจัยที่มีผลต่อคุณภาพของพืชหมักและการยับยั้งเชื้อสปอยเลจได้
    \item อธิบายกลไกของโพรไบโอติกในการปรับสมดุลจุลินทรีย์ในทางเดินอาหารสัตว์เคี้ยวเอื้องได้
    \item วางแผนกระบวนการเพาะเห็ดและการคัดเลือกวัสดุเพาะจากขี้เลื่อยไม้ยางพาราได้อย่างเหมาะสม
\end{enumerate}

\noindent\textbf{กิจกรรมการเรียนการสอน}
\begin{enumerate}
    \item การบรรยายและศึกษากระบวนการหมักหญ้าเนเปียร์ในถังหมักแบบจำลอง
    \item ปฏิบัติการวัดค่า pH และกรดแลกติกในตัวอย่างพืชหมัก
    \item การศึกษาโครงสร้างเส้นใยเห็ดนางรมบนวุ้นอาหารและการเจริญบนขี้เลื่อย
\end{enumerate}

\noindent\textbf{สื่อการเรียนการสอน}
\begin{enumerate}
    \item เอกสารคำสอนบทที่ 7 และแผนภูมิกระบวนการหมักเชิงชีวเคมี
    \item ตัวอย่างพืชหมักหญ้าเนเปียร์คุณภาพดีและพืชหมักที่เน่าเสีย
\end{enumerate}

\noindent\textbf{การประเมินผล}
\begin{enumerate}
    \item การประเมินผลการตรวจวิเคราะห์คุณภาพทางเคมีและชีววิทยาของพืชหมัก
    \item การทดสอบความรู้ความเข้าใจเกี่ยวกับการประยุกต์ใช้โพรไบโอติก
\end{enumerate}

\newpage

% ==========================================================
% เนื้อหาบทเรียน
% ==========================================================

\section{จุลชีววิทยาของการหมักพืชอาหารสัตว์}
\index{พืชอาหารสัตว์หมัก (Silage Fermentation)}

\textbf{พืชอาหารสัตว์หมัก (Silage)} เป็นวิธีการถนอมรักษาคุณค่าทางโภชนาการของพืชอาหารสัตว์ที่มีความชื้นสูง (เช่น ข้าวโพด หญ้าเนเปียร์ และยอดอ้อย) ไว้สำหรับเลี้ยงสัตว์เคี้ยวเอื้อง (โคเนื้อ โคนม) ในช่วงฤดูแล้ง โดยอาศัยการหมักของ \textbf{แบคทีเรียกรดแลกติก (Lactic Acid Bacteria: LAB)} ภายใต้สภาวะปิดที่ไร้ออกซิเจน (Anaerobic Condition)

เมื่อนำพืชอาหารสัตว์มาบดสับและอัดแน่นในหลุมหมัก (Silo) ออกซิเจนที่ค้างอยู่จะถูกใช้หมดไปอย่างรวดเร็วโดยการหายใจของเซลล์พืชและจุลินทรีย์กลุ่มแอโรบิก จากนั้นแบคทีเรียกรดแลกติก เช่น \textit{Lactobacillus plantarum}, \textit{Pediococcus acidilactici} และ \textit{Enterococcus faecium} จะใช้น้ำตาลละลายน้ำ (Water-Soluble Carbohydrates: WSC) ในเนื้อเยื่อพืชเป็นสารตั้งต้นในการสังเคราะห์กรดแลกติก (Lactic acid)

การสะสมของกรดแลกติกจะส่งผลให้ค่า pH ของพืชหมักลดลงอย่างรวดเร็วจาก 6.0 สู่ระดับต่ำกว่า 4.0--4.2 ซึ่งความเป็นกรดระดับนี้จะยับยั้งการเจริญเติบโตของแบคทีเรียเน่าเสีย สปอร์ของเชื้อ \textit{Clostridium butyricum} และเชื้อรา ทำให้พืชอาหารสัตว์สามารถเก็บรักษาไว้ได้นานนับปีโดยไม่สูญเสียคุณค่าทางอาหาร

\begin{figure}[H]
\centering
\begin{tikzpicture}[scale=0.95]
    % Sequence boxes
    \node (s1) [rectangle, rounded corners=6pt, draw=agriGreen, fill=agriMint, line width=1.2pt, minimum width=4.0cm, minimum height=1.2cm, align=center, font=\footnotesize]
        {\textbf{1. ระยะแอโรบิก (Aerobic)}\\เซลล์พืชหายใจ ใช้ออกซิเจนหมด\\อุณหภูมิสูงขึ้นเล็กน้อย (วัน 1--2)};

    \node (s2) [rectangle, rounded corners=6pt, draw=agriGold, fill=yellow!20, line width=1.2pt, minimum width=4.0cm, minimum height=1.2cm, align=center, font=\footnotesize, right=1.5cm of s1]
        {\textbf{2. ระยะหมักกรด (Fermentation)}\\แบคทีเรีย LAB สร้างกรดแลกติก\\pH ลดฮวบสู่ $< 4.0$ (วัน 3--14)};

    \node (s3) [rectangle, rounded corners=6pt, draw=agriLeaf, fill=agriLeaf!20, line width=1.2pt, minimum width=4.0cm, minimum height=1.2cm, align=center, font=\footnotesize, right=1.5cm of s2]
        {\textbf{3. ระยะคงตัว (Stable Phase)}\\เชื้อเน่าเสียถูกยับยั้งสมบูรณ์\\เก็บรักษาได้นานนับปี ($>14$ วัน)};

    % Arrows
    \draw[-{Stealth[scale=1.4]}, line width=2pt, color=agriGreen] (s1) -- (s2);
    \draw[-{Stealth[scale=1.4]}, line width=2pt, color=agriLeaf] (s2) -- (s3);

    % Graph of pH drop below
    \draw[-{Stealth[scale=1.1]}, line width=1.2pt, color=agriSlate] (0,-1.8) -- (14.0,-1.8) node[right, font=\tiny\bfseries] {เวลา (วัน)};
    \draw[-{Stealth[scale=1.1]}, line width=1.2pt, color=agriSlate] (0,-1.8) -- (0,-4.5) node[above, font=\tiny\bfseries] {ค่า pH};

    \node at (-0.4, -2.1) [font=\tiny] {6.5};
    \node at (-0.4, -3.2) [font=\tiny] {4.0};
    \node at (-0.4, -4.2) [font=\tiny] {3.5};

    % pH curve dropping rapidly
    \draw[line width=2pt, color=agriRed!90!black] 
        plot [smooth] coordinates {(0.2, -2.2) (2.0, -2.5) (4.5, -3.8) (7.0, -4.1) (13.5, -4.2)};
    \node at (6.5, -3.4) [font=\tiny\bfseries\color{agriRed!90!black}] {เส้นโค้งการลดลงของ pH};
    \draw[dashed, color=agriBlue] (0, -3.2) -- (13.5, -3.2);
    \node at (10.5, -3.0) [font=\tiny\color{agriBlue}] {ระดับปลอดภัยจาก Clostridium (pH $< 4.0$)};
\end{tikzpicture}
\caption{กระบวนการหมักพืชอาหารสัตว์ (Silage Fermentation) และพลวัตการลดลงของค่าความเป็นกรด-ด่าง}
\label{fig:silage_fermentation}
\end{figure}

\section{โพรไบโอติกในอาหารสัตว์}
\index{โพรไบโอติกในสัตว์ (Probiotics)}

หลังจากการประกาศแบนการใช้ยาปฏิชีวนะเพื่อเร่งการเจริญเติบโตในอาหารสัตว์ (Antibiotic Growth Promoters: AGPs) ในระดับสากล จุลินทรีย์โพรไบโอติก (Probiotics) ได้กลายมาเป็นทางเลือกทดแทนอันดับหนึ่ง สารชีวภัณฑ์โพรไบโอติกประกอบด้วยแบคทีเรียสายพันธุ์คัดเลือก เช่น \textit{Lactobacillus acidophilus}, \textit{Enterococcus faecium}, \textit{Bacillus subtilis} และยีสต์ \textit{Saccharomyces cerevisiae} ซึ่งมีบทบาทสำคัญ ดังนี้
\begin{itemize}
    \item การแข่งขันครอบครองพื้นที่ผิวเยื่อบุลำไส้ (Competitive Exclusion) กีดกันเชื้อก่อโรคในสัตว์ เช่น \textit{Salmonella} และ \textit{Escherichia coli}
    \item การผลิตกรดไขมันสายสั้น (Short-Chain Fatty Acids: SCFAs) ช่วยเพิ่มประสิทธิภาพการดูดซึมสารอาหาร
    \item การกระตุ้นระบบภูมิคุ้มกันประจำเยื่อเมือก (Mucosal Immunity) ในทางเดินอาหาร
\end{itemize}

\section{การเพาะเห็ดเศรษฐกิจกับการย่อยสลายลิกโนเซลลูโลส}
\index{การเพาะเห็ด (Mushroom Cultivation)}

เชื้อเห็ดจัดอยู่ในกลุ่มราชั้นสูง (Basidiomycota) มีความสามารถพิเศษในการสร้างเอนไซม์กลุ่มลิกโนไลติก (Lignolytic enzymes) ได้แก่ เอนไซม์แลกเคส (Laccase), ลิกนินเปอร์ออกซิเดส (Lignin Peroxidase) และแมงกานีสเปอร์ออกซิเดส (Manganese Peroxidase) ซึ่งสามารถตัดพันธะที่แข็งแรงในโมเลกุลลิกนินของขี้เลื่อยไม้ยางพารา ฟางข้าว และเปลือกข้าวโพด เปลี่ยนวัสดุเหลือทิ้งทางการเกษตรให้กลายเป็นดอกเห็ดที่มีคุณค่าทางโภชนาการสูง อุดมด้วยโปรตีน วิตามิน และสารเบตากลูแคน (Beta-glucan) เสริมภูมิคุ้มกัน

\begin{agribox}{การนำก้อนเชื้อเห็ดเก่ามาผลิตปุ๋ยหมักชีวภาพ}
ก้อนเชื้อเห็ดที่หมดอายุการให้ผลผลิต (Spent Mushroom Substrate: SMS) ยังคงมีเส้นใยราและเอนไซม์ย่อยสลายหลงเหลืออยู่เป็นจำนวนมาก เมื่อนำมาคลุกเคล้ากับมูลไก่และกากถั่วเหลือง แล้วผ่านการหมักแบบใช้ออกซิเจน จะได้ปุ๋ยอินทรีย์ที่มีโครงสร้างโปร่ง ร่วนซุย มีปริมาณอินทรียวัตถุสูง และมีเชื้อราปฏิปักษ์ช่วยป้องกันโรครากเน่าโคนเน่าได้อย่างยอดเยี่ยม
\end{agribox}

\section*{คำถามทบทวนท้ายบทที่ 7}
\addcontentsline{toc}{section}{คำถามทบทวนท้ายบทที่ 7}
\begin{enumerate}
    \item จงอธิบายบทบาทของแบคทีเรียกรดแลกติก (LAB) ในการถนอมอาหารสัตว์ และเหตุใดจึงต้องรักษาสภาพไร้ออกซิเจนในหลุมหมัก
    \item หากค่า pH ของพืชหมักไม่สามารถลดลงต่ำกว่า 4.5 ภายในระยะเวลา 7 วันแรก จะเกิดผลกระทบต่อคุณภาพของพืชหมักอย่างไร
    \item จงอธิบายกลไกของเอนไซม์กลุ่ม Lignolytic Enzymes ในเชื้อเห็ดต่อการย่อยสลายเศษวัสดุขี้เลื่อยไม้ยางพารา
\end{enumerate}

\clearpage